% =====================================================================
%  2bat-ud2-13.tex  —  SESSIÓ 13: Cas real complet amb GeoGebra
% =====================================================================
\horatitol{Sessió 13 — Cas real complet amb GeoGebra}%
{Integrar tot el mètode en un estudi de cas: de l'enunciat a la recomanació justificada, amb suport de GeoGebra.}

\subsection{Idea clau}

Avui tanquem el cicle complet d'una decisió basada en dades: de l'enunciat real a la
recomanació justificada. Fem \textbf{tot el procés} amb el suport de GeoGebra, que
ens estalvia errors de dibuix i ens deixa concentrar en el \textbf{plantejament} i la
\textbf{interpretació}.

\begin{idea}
\textbf{El cicle complet de la programació lineal}
\begin{enumerate}[leftmargin=1.6em, itemsep=2pt]
  \item \textbf{Variables:} definir què és $x$ i què és $y$.
  \item \textbf{Sistema:} traduir les condicions a inequacions (i la no-negativitat).
  \item \textbf{Regió factible:} representar-la (a GeoGebra) i \textbf{localitzar els
        vèrtexs} amb l'eina Intersecció.
  \item \textbf{Funció objectiu:} definir què s'optimitza i \textbf{avaluar-la als
        vèrtexs}.
  \item \textbf{Decisió:} triar l'òptim i \textbf{justificar} la recomanació.
\end{enumerate}
\end{idea}

\subsection{Exemples}

\begin{exemple}[un cas real, de cap a cap]
\textbf{Enunciat.} Un ajuntament reparteix places entre dos serveis: \textbf{atenció
domiciliària} ($x$) i \textbf{centre de dia} ($y$). El pressupost permet
$2x+3y\le 18$ (en milers d'euros) i el personal, $x+y\le 7$ places en total. Cada
plaça domiciliària atén $3$ persones i cada plaça de centre, $4$. Volem atendre el
màxim de persones: $P=3x+4y$.

\textbf{Regió factible i vèrtexs} (amb GeoGebra, eina Intersecció):

\begin{center}
\begin{tikzpicture}
\begin{axis}[udaxis, width=8.6cm, height=5.8cm,
    xlabel={\small domiciliària ($x$)}, ylabel={\small centre ($y$)},
    xmin=0, xmax=10, ymin=0, ymax=8, xtick={0,2,4,6,8}, ytick={0,2,4,6}]
  \fill[udblue!18] (axis cs:0,0) -- (axis cs:7,0) -- (axis cs:3,4) -- (axis cs:0,6) -- cycle;
  \addplot[udblue, domain=0:9, samples=2]{(18-2*x)/3};   % pressupost
  \addplot[udnavy, domain=0:7, samples=2]{7-x};          % personal
  \addplot[udnavy, only marks, mark=*, mark size=1.5pt] coordinates {(0,0)(7,0)(3,4)(0,6)};
  \node[font=\scriptsize, anchor=north] at (axis cs:7,-0.1) {$(7,0)$};
  \node[font=\scriptsize, anchor=south west] at (axis cs:3,4) {$(3,4)$};
  \node[font=\scriptsize, anchor=east] at (axis cs:0,6) {$(0,6)$};
\end{axis}
\end{tikzpicture}
\end{center}

\textbf{Avaluem} $P=3x+4y$ als vèrtexs:
\[
P(0,0)=0,\quad P(7,0)=21,\quad P(3,4)=9+16=25,\quad P(0,6)=24.
\]
El màxim és $P=25$ al vèrtex $(3,4)$.
\end{exemple}

\begin{exemple}[la recomanació justificada]
\textbf{Decisió:} obrir \textbf{$3$ places domiciliàries i $4$ de centre de dia}, que
atenen \textbf{$25$ persones}, el màxim possible. \textbf{Justificació:} aquesta
combinació esgota tots dos recursos (pressupost $2\cdot 3+3\cdot 4=18$ i personal
$3+4=7$), i supera les alternatives extremes (només domiciliària, $21$; només centre,
$24$). Equilibrar els dos serveis atén més gent que concentrar-ho en un de sol.
\end{exemple}

\subsection{Exercicis resolts}

\begin{exresolt}[verificar l'òptim a GeoGebra]
Com comprovaries a GeoGebra que el vèrtex $(3,4)$ és correcte i que és l'òptim?
\solucio
A GeoGebra: (1) introdueixes les inequacions \texttt{2x+3y<=18}, \texttt{x+y<=7},
\texttt{x>=0}, \texttt{y>=0} i veus la regió; (2) amb l'eina \emph{Intersecció} sobre
les rectes \texttt{2x+3y=18} i \texttt{x+y=7} obtens el punt $(3,4)$; (3) calcules
$P=3x+4y$ a cada vèrtex (pots escriure \texttt{P(3,4)} si has definit la funció) i
comproves que $25$ és el valor més gran. La màquina confirma el càlcul fet a mà.
\resposta{Amb Intersecció es localitza $(3,4)$ i, avaluant $P$ als vèrtexs, es confirma que és el màxim.}
\end{exresolt}

\begin{exresolt}[i si canvia un recurs?]
Si el personal permetés $x+y\le 8$ (una plaça més), com afectaria l'òptim? (Recalcula
el vèrtex de creuament.)
\solucio
La restricció de personal puja a $x+y=8$. Creuament amb $2x+3y=18$: de $y=8-x$,
$2x+3(8-x)=18\Rightarrow -x+24=18\Rightarrow x=6$, $y=2$, vèrtex $(6,2)$. Avaluem
$P=3x+4y$ als nous vèrtexs $(0,0)$, $(8,0)$, $(6,2)$, $(0,6)$:
$P(8,0)=24$, $P(6,2)=18+8=26$, $P(0,6)=24$. Ara el màxim és $P=26$ a $(6,2)$: una
plaça més de personal permet atendre una persona més.
\resposta{Amb $x+y\le 8$, l'òptim passa a $(6,2)$ amb $P=26$.}
\end{exresolt}

\subsection{Exercicis per fer}

\noindent\textit{Cas: $2x+3y\le 18$, $x+y\le 7$, $x,y\ge 0$; vèrtexs $(0,0)$, $(7,0)$,
$(3,4)$, $(0,6)$.}

\begin{exercicis}
\begin{exlist}
  \item Comprova que els quatre vèrtexs compleixen totes les restriccions (substitueix
        a les dues inequacions).
  \item Avalua $P=3x+4y$ als quatre vèrtexs i confirma que el màxim és a $(3,4)$.
  \item Suposa ara que cada plaça domiciliària atén $5$ persones i cada plaça de centre
        $4$ ($P=5x+4y$). Troba el nou òptim. Ha canviat?
  \item Si en lloc de maximitzar persones volguéssim \textbf{minimitzar} un cost
        $C=2x+3y$ sobre la mateixa regió, on seria el mínim? (Compte: és l'altre
        objectiu.)
  \item \textbf{(Recomanació)} Per a l'òptim de l'exercici 2, redacta una recomanació
        breu i justificada per a l'ajuntament (quantes places de cada servei, quantes
        persones ateses, quins recursos s'esgoten).
  \item \textbf{(Pràctica GeoGebra)} Descriu els passos que faries a GeoGebra per
        resoldre aquest cas: quines ordres escriuries i quina eina faries servir per
        trobar els vèrtexs.
\end{exlist}
\end{exercicis}

% ---------------------------------------------------------------------
\fullsolucions{Sessió 13}
\begin{solucions}
\begin{sollist}
  \item $(0,0)$: $0\le 18$, $0\le 7$. $(7,0)$: $14\le 18$, $7\le 7$. $(3,4)$:
        $18\le 18$, $7\le 7$. $(0,6)$: $18\le 18$, $6\le 7$. Tots compleixen.
  \item $P(0,0)=0$, $P(7,0)=21$, $P(3,4)=25$, $P(0,6)=24$. Màxim $P=25$ a $(3,4)$.
  \item $P=5x+4y$: $P(7,0)=35$, $P(3,4)=15+16=31$, $P(0,6)=24$. Ara el màxim és
        $P=35$ a $(7,0)$: com que la domiciliària atén més per plaça, convé concentrar-hi
        els recursos. L'òptim \textbf{ha canviat} a $(7,0)$.
  \item $C=2x+3y$: $C(0,0)=0$, $C(7,0)=14$, $C(3,4)=18$, $C(0,6)=18$. Mínim $C=0$ a
        $(0,0)$ (no obrir places); si s'exclou el trivial, el menor és $C=14$ a $(7,0)$.
  \item Resposta oberta. Ha d'indicar $3$ places domiciliàries i $4$ de centre, $25$
        persones ateses, i que s'esgoten pressupost ($18$) i personal ($7$).
  \item Resposta oberta. Passos esperats: introduir \texttt{2x+3y<=18}, \texttt{x+y<=7},
        \texttt{x>=0}, \texttt{y>=0}; usar l'eina \emph{Intersecció} sobre les rectes
        per als vèrtexs; definir \texttt{P(x,y)=3x+4y} i avaluar-la a cada vèrtex.
\end{sollist}
\end{solucions}
